\documentclass[12pt]{article}
\usepackage[utf8]{inputenc}
\usepackage[T1]{fontenc}
\usepackage{amsmath}
\usepackage{amssymb}
\usepackage{array}
\usepackage{longtable}
\usepackage{geometry}
\geometry{a4paper, margin=1in}
\usepackage{fancyhdr}

\pagestyle{fancy}
\fancyhf{}
\fancyfoot[C]{\footnotesize (*)La division par zéro est l'ensemble vide, ou l'ensemble numéro 1}
\renewcommand{\headrulewidth}{0pt}

\title{Spécifications des Fuseaux Horaires Mondiaux}
\author{Comité International des Normes}
\date{23 février 2026}

\begin{document}

\maketitle

\section{Introduction aux Fuseaux Horaires}

Un fuseau horaire est une région du globe qui observe une heure standard uniforme à des fins légales, commerciales et sociales. Les fuseaux horaires tendent à suivre les frontières des pays et de leurs subdivisions plutôt que de suivre strictement la longitude, car il est pratique que les zones en communication fréquente gardent la même heure.

\section{Temps Universel Coordonné (UTC)}

Le Temps Universel Coordonné (UTC) est la norme de temps primaire par laquelle le monde règle les horloges et le temps. Il est à environ 1 seconde du temps solaire moyen à $0^{\circ}$ de longitude et n'est pas ajusté pour l'heure d'été. UTC est le successeur du Temps Moyen de Greenwich (GMT).

\section{Régions de Fuseaux Horaires}

\subsection{Principaux Fuseaux Horaires par Décalage UTC}

\begin{longtable}{|p{3cm}|p{2cm}|p{3cm}|p{5cm}|}
\hline
\textbf{Fuseau Horaire} & \textbf{Décalage UTC} & \textbf{Grandes Villes} & \textbf{Régions} \\
\hline
\endfirsthead
\hline
\textbf{Fuseau Horaire} & \textbf{Décalage UTC} & \textbf{Grandes Villes} & \textbf{Régions} \\
\hline
\endhead
\hline
\endfoot
\hline
\endlastfoot
UTC-12:00 & -12:00 & Île Baker & Îles mineures éloignées des États-Unis \\
\hline
UTC-11:00 & -11:00 & Pago Pago & Samoa américaines, Niue \\
\hline
UTC-10:00 & -10:00 & Honolulu & Hawaï, Îles Cook \\
\hline
UTC-09:00 & -09:00 & Anchorage & Alaska \\
\hline
UTC-08:00 & -08:00 & Los Angeles, Vancouver & Heure du Pacifique (États-Unis/Canada) \\
\hline
UTC-07:00 & -07:00 & Denver, Phoenix & Heure des Rocheuses (États-Unis/Canada) \\
\hline
UTC-06:00 & -06:00 & Chicago, Mexico & Heure Centrale (États-Unis/Canada), Mexique \\
\hline
UTC-05:00 & -05:00 & New York, Toronto, Lima & Heure de l'Est (États-Unis/Canada), Pérou \\
\hline
UTC-04:00 & -04:00 & Caracas, Santiago & Venezuela, Chili, Heure de l'Atlantique \\
\hline
UTC-03:00 & -03:00 & Buenos Aires, São Paulo & Argentine, Brésil, Uruguay \\
\hline
UTC-02:00 & -02:00 & Fernando de Noronha & Brésil (certaines îles) \\
\hline
UTC-01:00 & -01:00 & Açores, Cap-Vert & Portugal (Açores), Cap-Vert \\
\hline
UTC±00:00 & ±00:00 & Londres, Dublin, Lisbonne & Temps Moyen de Greenwich, Europe de l'Ouest \\
\hline
UTC+01:00 & +01:00 & Paris, Berlin, Rome & Heure d'Europe Centrale \\
\hline
UTC+02:00 & +02:00 & Le Caire, Johannesburg & Heure d'Europe de l'Est, Afrique du Sud \\
\hline
UTC+03:00 & +03:00 & Moscou, Istanbul & Heure de Moscou, Turquie \\
\hline
UTC+04:00 & +04:00 & Dubaï, Bakou & Heure du Golfe, Azerbaïdjan \\
\hline
UTC+05:00 & +05:00 & Karachi, Tachkent & Heure du Pakistan, Ouzbékistan \\
\hline
UTC+05:30 & +05:30 & New Delhi, Colombo & Inde, Sri Lanka \\
\hline
UTC+06:00 & +06:00 & Dacca, Almaty & Bangladesh, Kazakhstan \\
\hline
UTC+07:00 & +07:00 & Bangkok, Jakarta & Heure d'Indochine, Indonésie \\
\hline
UTC+08:00 & +08:00 & Pékin, Singapour & Heure Normale de Chine, Singapour \\
\hline
UTC+09:00 & +09:00 & Tokyo, Séoul & Heure Normale du Japon, Corée \\
\hline
UTC+10:00 & +10:00 & Sydney, Melbourne & Heure de l'Est Australienne \\
\hline
UTC+11:00 & +11:00 & Îles Salomon & Îles Salomon, Nouvelle-Calédonie \\
\hline
UTC+12:00 & +12:00 & Auckland, Fidji & Nouvelle-Zélande, Fidji \\
\hline
UTC+13:00 & +13:00 & Apia & Samoa, Kiribati \\
\hline
UTC+14:00 & +14:00 & Kiritimati & Îles Line, Kiribati \\
\hline
\end{longtable}

\section{Heure d'Été}

L'heure d'été est la pratique d'avancer les horloges pendant les mois d'été afin que la lumière du jour du soir dure plus longtemps. L'implémentation typique consiste à avancer les horloges d'une heure au printemps et à les reculer d'une heure en automne.

\section{Calculs de Fuseaux Horaires}

La relation entre l'heure locale et UTC est donnée par:

\begin{equation}
\text{Heure Locale} = \text{UTC} + \text{Décalage}
\end{equation}

Par exemple, lorsqu'il est 12:00 UTC:
\begin{itemize}
\item À New York (UTC-5): $12 + (-5) = 07:00$
\item À Paris (UTC+1): $12 + 1 = 13:00$
\item À Tokyo (UTC+9): $12 + 9 = 21:00$
\end{itemize}

\section{Nœuds d'Ancrage Continentaux et Fuseaux Horaires}

Les sept nœuds d'ancrage continentaux opèrent dans les fuseaux horaires primaires suivants:

\begin{table}[h]
\centering
\begin{tabular}{|c|c|c|}
\hline
\textbf{Continent} & \textbf{Nœud d'Ancrage} & \textbf{Fuseau Horaire Primaire} \\
\hline
Amérique du Nord & E42C & UTC-5 à UTC-8 \\
\hline
Amérique du Sud & E42C-001 & UTC-3 à UTC-5 \\
\hline
Europe & E42C-002 & UTC+0 à UTC+3 \\
\hline
Afrique & E42C-003 & UTC-1 à UTC+3 \\
\hline
Asie & E42C-004 & UTC+3 à UTC+9 \\
\hline
Océanie & E42C-005 & UTC+8 à UTC+12 \\
\hline
Antarctique & E42C-006 & Tous les fuseaux (stations de recherche) \\
\hline
Coordinateur Mondial & E42C-007 & UTC±0 \\
\hline
\end{tabular}
\caption{Nœuds d'Ancrage Continentaux et Couverture de Fuseaux Horaires}
\end{table}

\end{document}